\chapter{Verificação de tipos}

\section{Introdução}

No capítulo anterior, definimos formalmente o que é um sistema de
tipos e demonstramos as propriedades de corretude que ele oferece.
Vimos que as regras de tipagem, escritas como sistemas de inferência,
constituem uma especificação matemática precisa do que significa um
programa estar bem tipado. Este capítulo trata de transformar essa
especificação em código: o \textbf{verificador de tipos} (\emph{type checker}).

A ideia central é que existe uma correspondência direta entre as
regras de inferência e as equações de um programa funcional. Cada
regra se torna um caso em uma função recursiva; as premissas da regra
se tornam chamadas recursivas; e a conclusão é o resultado retornado.
Essa tradução é tão sistemática que, em muitos casos, o verificador
de tipos pode ser lido como uma transcrição quase literal das regras.

Apresentamos verificadores de tipos para quatro linguagens em
progressão crescente de complexidade: \emph{TExp}, \emph{TLine},
\emph{TWhile} e \emph{TImp}. Os sistemas de tipos formais dessas
linguagens foram apresentados no capítulo anterior; aqui focamos
exclusivamente na tradução dessas regras para código Haskell.

\section{Verificação de Tipos para TExp}

\subsection{A Linguagem e o seu Sistema de Tipos}

A linguagem TExp tem dois tipos, \(\mathbf{Bool}\) e
\(\mathbf{Nat}\), e sete construções: \(\mathbf{true}\),
\(\mathbf{false}\), \(\mathbf{if}\), \(\mathbf{0}\),
\(\mathbf{succ}\), \(\mathbf{pred}\) e \(\mathbf{iszero}\). Em Haskell:

\begin{minted}{haskell}
data Ty   = TBool | TNat

data Term
  = TTrue  | TFalse
  | TIf Term Term Term
  | TZero
  | TSucc Term | TPred Term | TIsZero Term
\end{minted}

O sistema de tipos é definido pela relação \(\vdash t : T\) através
de sete regras de inferência. A verificação de tipos é o problema de,
dado um termo \(t\), encontrar \(T\) tal que \(\vdash t : T\), ou
reportar que nenhum tal \(T\) existe.

\subsection{Da Regra à Equação}

A tradução de cada regra de tipagem para uma equação Haskell segue um
padrão uniforme. Consideremos a regra (T-If):

\[
  \frac{
    \vdash t_1 : \mathbf{Bool} \quad
    \vdash t_2 : T \quad
    \vdash t_3 : T
  }{
    \vdash \mathbf{if}\ t_1\ \mathbf{then}\ t_2\ \mathbf{else}\ t_3 : T
  } \quad\text{(T-If)}
\]

A conclusão nos diz que o padrão a combinar é
\haskell{TIf t1 t2 t3}. As premissas indicam três
verificações recursivas a realizar:

\begin{enumerate}
  \item Verificar que \(t_1\) tem tipo \(\mathbf{Bool}\): a condição
    deve ser booleana;
  \item Verificar que \(t_2\) tem algum tipo \(T\);
  \item Verificar que \(t_3\) tem o \textbf{mesmo} tipo \(T\): os
    ramos devem coincidir;
  \item Retornar \(T\) como resultado.
\end{enumerate}

Em Haskell:

\begin{minted}{haskell}
tcTerm (TIf t1 t2 t3) = do
    ty1 <- tcTerm t1
    unless (ty1 == TBool) $
        throwError "Type error in if: condition must have type Bool"
    ty2 <- tcTerm t2
    ty3 <- tcTerm t3
    unless (ty2 == ty3) $
        throwError "Type error in if: both branches must have the same type"
    pure ty2
\end{minted}

A mônada \haskell{TcM = ExceptT String Identity}
carrega a possibilidade de falha sem precisar de nenhum contexto
externo, pois TExp não tem variáveis.

\subsection{O Verificador Completo}

As demais equações seguem o mesmo padrão: cada uma implementa
diretamente a sua regra correspondente:

\begin{minted}{haskell}
type TcM a = ExceptT String Identity a

tcTerm :: Term -> TcM Ty
tcTerm TTrue  = pure TBool                          -- (T-True)
tcTerm TFalse = pure TBool                          -- (T-False)
tcTerm TZero  = pure TNat                           -- (T-Zero)

tcTerm (TIf t1 t2 t3) = do                         -- (T-If)
    ty1 <- tcTerm t1
    unless (ty1 == TBool) $
        throwError "Type error in if: condition must have type Bool"
    ty2 <- tcTerm t2
    ty3 <- tcTerm t3
    unless (ty2 == ty3) $
        throwError "Type error in if: both branches must have the same type"
    pure ty2

tcTerm (TSucc t1) = do                              -- (T-Succ)
    ty <- tcTerm t1
    unless (ty == TNat) $
        throwError "Type error in succ: argument must have type Nat"
    pure TNat

tcTerm (TPred t1) = do                              -- (T-Pred)
    ty <- tcTerm t1
    unless (ty == TNat) $
        throwError "Type error in pred: argument must have type Nat"
    pure TNat

tcTerm (TIsZero t1) = do                            -- (T-IsZero)
    ty <- tcTerm t1
    unless (ty == TNat) $
        throwError "Type error in iszero: argument must have type Nat"
    pure TBool
\end{minted}

A estrutura recursiva de \haskell{tcTerm} espelha a
estrutura indutiva do sistema de tipos. O verificador de tipos é, em
essência, a prova construtiva do lema de unicidade de tipos,
executada como um programa.

\section{Verificação de Tipos para TLine}

\subsection{A Linguagem TLine}

TLine é uma linguagem imperativa simples com declaração de variáveis
tipadas, atribuição, leitura e impressão. A sua sintaxe abstrata é:

\begin{minted}{haskell}
data Ty   = TInt | TBool | TString

data Stmt
  = SDecl   Var Ty Exp   -- var x : T = e ;
  | SAssign Var Exp      -- x := e ;
  | SRead   Exp Var      -- read prompt x ;
  | SPrint  Exp          -- print e ;

data Exp
  = EInt Int | EBool Bool | EString String
  | EVar Var
  | Exp :+: Exp | Exp :*: Exp | Exp :-: Exp | Exp :/: Exp
  | Exp :<: Exp | Exp :=: Exp
  | Not Exp
\end{minted}

\subsection{O Contexto de Tipos}

A principal novidade em relação a TExp é a presença de
\textbf{variáveis}. Para verificar o tipo de
\haskell{EVar x}, é preciso saber o tipo declarado
para \haskell{x}. Essa informação é mantida em um
\textbf{contexto de tipos}:

\[
  \Gamma : \mathit{Var} \rightharpoonup \mathit{Ty}
\]

Em Haskell, o contexto é representado por um \haskell{Map}:

\begin{minted}{haskell}
type Ctx = Map Var Ty
\end{minted}

O verificador precisa tanto \textbf{consultar} o contexto (para
variáveis) quanto \textbf{estendê-lo} (quando uma declaração é
processada). A mônada combina estado e tratamento de erros:

\begin{minted}{haskell}
type TcM a = StateT Ctx (ExceptT String Identity) a
\end{minted}

O \haskell{StateT Ctx} transporta o contexto ao longo
da verificação; o \haskell{ExceptT String} interrompe
a verificação com uma mensagem de erro.

\subsection{O Algoritmo de Verificação}

O verificador implementa diretamente as regras. Para expressões:

\begin{minted}{haskell}
tcExp :: Exp -> TcM Ty
tcExp (EInt _)    = pure TInt        -- (T-Int)
tcExp (EBool _)   = pure TBool       -- (T-Bool)
tcExp (EString _) = pure TString    -- (T-String)
tcExp (EVar v)    = lookupVar v    -- (T-Var)
tcExp (e1 :+: e2) = do             -- (T-ArithOp)
    ty1 <- tcExp e1
    ty2 <- tcExp e2
    unless (ty1 == TInt) $ throwError "Type error"
    unless (ty2 == TInt) $ throwError "Type error"
    pure TInt
tcExp (e1 :<: e2) = do                                -- (T-CmpOp)
    ty1 <- tcExp e1
    ty2 <- tcExp e2
    unless (ty1 == TInt) $ throwError "Type error"
    unless (ty2 == TInt) $ throwError "Type error"
    pure TBool
tcExp (Not e1) = do                                   -- (T-Not)
    ty <- tcExp e1
    unless (ty == TBool) $ throwError "Type error"
    pure TBool
\end{minted}

As equações para \haskell{:*:},
\haskell{:-:} e \haskell{:/:} são
idênticas à de \haskell{:+:}; a de
\haskell{:=:} é idêntica à de \haskell{:<:}.

Para comandos, o contexto é atualizado com
\haskell{modify} quando uma variável é declarada:

\begin{minted}{haskell}
tcStmt :: Stmt -> TcM ()
tcStmt (SDecl v t e) = do                             -- (T-Decl)
    t' <- tcExp e
    unless (t == t') $ throwError "Type error"
    addNewVar v t

tcStmt (SAssign v e) = do                             -- (T-Assign)
    t  <- lookupVar v
    t' <- tcExp e
    unless (t == t') $ throwError "Type error"

tcStmt (SRead e v) = do                               -- (T-Read)
    t <- tcExp e
    unless (t == TString) $ throwError "Type error"
    t' <- lookupVar v
    unless (t' == TInt)   $ throwError "Type error"

tcStmt (SPrint e) = tcExp e >> pure ()                -- (T-Print)
\end{minted}

O ponto de entrada processa os comandos sequencialmente e retorna o
contexto final:

\begin{minted}{haskell}
typeCheck :: TLine -> Either String [(Var, Ty)]
typeCheck p = either Left (Right . Map.toList) (runTcM (tcTLine p))

tcTLine :: TLine -> TcM ()
tcTLine (TLine ss) = mapM_ tcStmt ss

runTcM :: TcM a -> Either String Ctx
runTcM m = runIdentity (runExceptT (execStateT m Map.empty))
\end{minted}

A função \haskell{addNewVar} estende o contexto com o par \(x
\mapsto T\), correspondendo à notação \(\Gamma[x \mapsto T]\) nas
regras. \haskell{lookupVar} consulta o contexto e
lança um erro se a variável não estiver definida, implementando a
premissa \(x \in \mathrm{dom}(\Gamma)\).

\section{Verificação de Tipos para TWhile}

\subsection{Extensões à Linguagem}

TWhile estende TLine com dois novos comandos:

\begin{minted}{haskell}
data Stmt
  = SDecl   Var Ty Exp
  | SAssign Var Exp
  | SRead   Exp Var
  | SPrint  Exp
  | SIf     Exp Block Block
  | SWhile  Exp Block

type Block = [Stmt]
\end{minted}

\subsection{Escopo de Variáveis}

A novidade mais importante de TWhile é que variáveis declaradas
\textbf{dentro} de um bloco \haskell{if} ou
\haskell{while} não devem ser visíveis fora dele.
Este é o conceito de \textbf{escopo léxico}: cada bloco tem o seu
próprio espaço de nomes local.

A regra se traduz na seguinte exigência de implementação: ao entrar
em um bloco, o contexto pode ser estendido pelas declarações
internas; ao sair do bloco, o contexto deve ser restaurado ao estado
anterior. Essa operação é implementada por
\haskell{withLocalCtx}:

\begin{minted}{haskell}
withLocalCtx :: TcM a -> TcM a
withLocalCtx m = do
    ctx <- get     -- salva o contexto atual
    r   <- m       -- verifica o bloco
    put ctx        -- restaura o contexto original
    pure r
\end{minted}

E \haskell{tcBlock} aplica essa operação a cada bloco
de comandos:

\begin{minted}{haskell}
tcBlock :: [Stmt] -> TcM ()
tcBlock ss = withLocalCtx (mapM_ tcStmt ss)
\end{minted}

\subsection{As Novas Equações do Verificador}

\begin{minted}{haskell}
tcStmt :: Stmt -> TcM ()
-- ... equações idênticas as de TLine ...
tcStmt (SIf e b1 b2) = do  -- (T-If)
    t <- tcExp e
    unless (t == TBool) $ throwError "Type error"
    tcBlock b1
    tcBlock b2

tcStmt (SWhile e b) = do   -- (T-While)
    t <- tcExp e
    unless (t == TBool) $ throwError "Type error"
    tcBlock b
\end{minted}

A correspondência com as regras é imediata:
\haskell{tcExp e} verifica a premissa \(\Gamma \vdash
e : \mathbf{Bool}\); as chamadas a \haskell{tcBlock}
verificam os blocos; e como \haskell{withLocalCtx}
restaura o contexto após cada bloco, o escopo léxico é automaticamente imposto.

% ============================================================
\section{Verificação de Tipos para TImp}
% ============================================================

TImp estende TWhile com declarações de funções e tipos registro. A
sua verificação de tipos introduz dois novos ambientes globais
(\(\Delta\) e \(\Phi\)) e um novo componente local (o tipo de retorno
esperado da função corrente \(\rho_{\mathit{ret}}\)). A implementação
encontra-se em \verb|TImp/Frontend/TypeCheck/TImpTypeCheck.hs|.

% ------------------------------------------------------------
\subsection{A Mônada de Verificação}
% ------------------------------------------------------------

TWhile usa \haskell{StateT Ctx (ExceptT String Identity)} para
transportar o contexto de variáveis mutável. TImp precisa, adicionalmente,
de ambientes \emph{imutáveis}: \(\Delta\), \(\Phi\) e
\(\rho_{\mathit{ret}}\) nunca mudam dentro de um bloco. A estrutura
natural para informação somente-leitura é \haskell{ReaderT}:

\begin{minted}{haskell}
data TcEnv = TcEnv
  { recordEnv  :: Map Name [(Field, Ty)]  -- Delta
  , funcEnv    :: Map Name ([Ty], RetTy)  -- Phi
  , returnType :: Maybe RetTy             -- rho_ret
  }

type VarCtx = Map Var Ty

type TcM a = ReaderT TcEnv
               (StateT VarCtx
                 (ExceptT String Identity)) a
\end{minted}

A pilha combina três camadas:

\begin{itemize}
  \item \haskell{ReaderT TcEnv}: ambientes globais imutáveis
    (\(\Delta\), \(\Phi\)) e o tipo de retorno esperado; acessados com
    \haskell{asks}.
  \item \haskell{StateT VarCtx}: contexto de variáveis \(\Gamma\),
    que cresce à medida que variáveis são declaradas; manipulado com
    \haskell{get}/\haskell{modify}.
  \item \haskell{ExceptT String}: erros de tipo, interrompidos com
    \haskell{throwError}.
\end{itemize}

O ponto de entrada constrói \(\Delta\) e \(\Phi\) a partir das
declarações antes de iniciar a verificação:

\begin{minted}{haskell}
typeCheck :: TImp -> Either String TcResult
typeCheck (TImp decls stmts) = do
  recEnv  <- buildRecordEnv decls      -- constroi Delta
  funcEnv <- buildFuncEnv recEnv decls -- constroi Phi
  let env = TcEnv recEnv funcEnv Nothing
  mapM_ (tcFuncDecl env) [fd | DFunc fd <- decls]
  finalCtx <- runTcM env (tcBlock stmts)
  ...
\end{minted}

A construção separada de \(\Delta\) e \(\Phi\) garante que
declarações de funções e registros sejam visíveis mutuamente:
qualquer função pode referenciar qualquer tipo registro e qualquer
outra função, independentemente da ordem de declaração.

% ------------------------------------------------------------
\subsection{Regras e Implementação: Registros}
% ------------------------------------------------------------

\subsubsection{Acesso a campo: regra (T-Field)}

\[
  \infer[\text{(T-Field)}]
  {\Gamma;\Delta;\Phi \vdash e.f : \tau}
  {\Gamma;\Delta;\Phi \vdash e : R \quad \Delta(R)(f) = \tau}
\]

\begin{minted}{haskell}
tcExp (EField e f) = do
  t <- tcExp e
  case t of
    TRecord rname -> lookupField rname f   -- Delta(R)(f)
    _ -> throwError ("Field access on non-record type " ++ showTy t)
\end{minted}

\haskell{lookupField} consulta \(\Delta\) via
\haskell{asks recordEnv} e lança um erro se o campo não existe.

\subsubsection{Construção de registro: regra (T-New)}

\[
  \infer[\text{(T-New)}]
  {\Gamma;\Delta;\Phi \vdash \mathbf{new}\;R\,\{\overline{f_i\!=\!e_i}\}: R}
  {R \in \mathrm{dom}(\Delta)
    \quad \Delta(R) = \overline{(f_i:\tau_i)}
  \quad \Gamma;\Delta;\Phi \vdash e_i : \tau_i}
\]

\begin{minted}{haskell}
tcExp (ENew rname fields) = do
  renv <- asks recordEnv
  case Map.lookup rname renv of
    Nothing -> throwError ("Undefined record type: " ++ rname)
    Just declFlds -> do
      mapM_ (checkField declFlds) fields  -- verifica tipos dos campos
      -- verifica que nenhum campo extra foi fornecido
      let provided = map fst fields
          declared = map fst declFlds
      unless (all (`elem` declared) provided) $
        throwError ("Unknown field in new " ++ rname)
      unless (length provided == length declared) $
        throwError ("Wrong number of fields in new " ++ rname)
      pure (TRecord rname)
\end{minted}

\subsubsection{Atribuição de campo: regra (T-FieldAssign)}

\[
  \infer[\text{(T-FieldAssign)}]
  {\Gamma;\Delta;\Phi \vdash x.f \mathbin{:=} e \dashv \Gamma}
  {\Gamma(x) = R \quad \Delta(R)(f) = \tau
  \quad \Gamma;\Delta;\Phi \vdash e : \tau}
\]

\begin{minted}{haskell}
tcStmt (SFieldAssign v f e) = do
  vTy <- lookupVar v
  case vTy of
    TRecord rname -> do
      fieldTy <- lookupField rname f     -- Delta(R)(f)
      t' <- tcExp e
      unless (fieldTy == t') $
        throwError ("Type error in field assignment " ++ v ++ "." ++ f)
    _ -> throwError (v ++ " is not a record")
\end{minted}

% ------------------------------------------------------------
\subsection{Regras e Implementação: Funções}
% ------------------------------------------------------------

\subsubsection{Chamada como expressão: regra (T-CallExpr)}

\[
  \infer[\text{(T-CallExpr)}]
  {\Gamma;\Delta;\Phi \vdash f(e_1,\ldots,e_n) : \tau}
  {\Phi(f) = ([\tau_1,\ldots,\tau_n],\;\tau)
  \quad \Gamma;\Delta;\Phi \vdash e_i : \tau_i}
\]

\begin{minted}{haskell}
tcExp (ECall fname args) = do
  fenv <- asks funcEnv
  case Map.lookup fname fenv of
    Nothing -> throwError ("Undefined function: " ++ fname)
    Just (pts, rt) -> do
      argTys <- mapM tcExp args
      unless (argTys == pts) $
        throwError ("Argument type mismatch in call to " ++ fname)
      case rt of
        RTVoid  -> throwError ("Void function " ++ fname ++ " used as expression")
        RTTy ty -> pure ty
\end{minted}

A verificação de que \(rt \neq \mathbf{void}\) implementa a distinção
entre (T-CallExpr) e (T-CallStmt): uma função \texttt{void} não pode
aparecer onde um valor é esperado.

\subsubsection{Chamada como comando: regra (T-CallStmt)}

\[
  \infer[\text{(T-CallStmt)}]
  {\Gamma;\Delta;\Phi \vdash f(e_1,\ldots,e_n) \dashv \Gamma}
  {\Phi(f) = ([\tau_1,\ldots,\tau_n],\;\mathbf{void})
  \quad \Gamma;\Delta;\Phi \vdash e_i : \tau_i}
\]

\begin{minted}{haskell}
tcStmt (SCall fname args) = do
  fenv <- asks funcEnv
  case Map.lookup fname fenv of
    Nothing -> throwError ("Undefined function: " ++ fname)
    Just (pts, rt) -> do
      unless (rt == RTVoid) $
        throwError ("Non-void call used as statement: " ++ fname)
      argTys <- mapM tcExp args
      unless (argTys == pts) $
        throwError ("Argument type mismatch in call to " ++ fname)
\end{minted}

\subsubsection{Retorno: regras (T-RetVal) e (T-RetVoid)}

\[
  \infer[\text{(T-RetVal)}]
  {\ldots;\rho_{\mathit{ret}} \vdash \mathbf{return}\; e \dashv \Gamma}
  {\rho_{\mathit{ret}} = \tau \quad \Gamma;\Delta;\Phi \vdash e : \tau}
  \qquad
  \infer[\text{(T-RetVoid)}]
  {\ldots;\mathbf{void} \vdash \mathbf{return} \dashv \Gamma}
  {}
\]

O tipo de retorno esperado é lido do ambiente via \haskell{asks
returnType}. O valor \haskell{Nothing} indica que estamos fora de
qualquer função; \haskell{return} fora de função é um erro:

\begin{minted}{haskell}
tcStmt (SReturn Nothing) = do
  mrt <- asks returnType
  case mrt of
    Nothing       -> throwError "return outside a function"
    Just RTVoid   -> pure ()
    Just (RTTy t) -> throwError ("return with no value in function returning "
                                  ++ showTy t)

tcStmt (SReturn (Just e)) = do
  mrt <- asks returnType
  t'  <- tcExp e
  case mrt of
    Nothing       -> throwError "return outside a function"
    Just RTVoid   -> throwError "return with value in void function"
    Just (RTTy t) ->
      unless (t == t') $
        throwError ("return type mismatch: expected " ++ showTy t)
\end{minted}

\subsubsection{Declaração de função: regra (T-FuncDecl)}

\[
  \infer[\text{(T-FuncDecl)}]
  {\Delta;\Phi \vdash \mathbf{fn}\; f(\overline{x_i:\tau_i}):\rho\;\{B\}
  \;\checkmark}
  {[x_1\!\mapsto\!\tau_1,\ldots,x_n\!\mapsto\!\tau_n];\Delta;\Phi;\rho
  \vdash B \dashv \_}
\]

A função \haskell{tcFuncDecl} instala o contexto inicial com os
parâmetros e ajusta \(\rho_{\mathit{ret}}\) no ambiente via
\haskell{local}:

\begin{minted}{haskell}
tcFuncDecl :: TcEnv -> FuncDecl -> Either String ()
tcFuncDecl baseEnv (FuncDecl name params retTy body) = do
  let paramPairs = [(v, t) | Param v t <- params]
      initCtx    = Map.fromList paramPairs
      env        = baseEnv { returnType = Just retTy }
  case runIdentity (runExceptT
         (execStateT (runReaderT (tcBlock body) env) initCtx)) of
    Left err -> Left ("In function " ++ name ++ ": " ++ err)
    Right _  -> Right ()
\end{minted}

O contexto de variáveis é inicializado com os parâmetros formais
(\(\Gamma_0 = [x_1\mapsto\tau_1,\ldots,x_n\mapsto\tau_n]\)) e o
\haskell{returnType} é fixado a \haskell{Just retTy}. Como o ambiente
é imutável e os parâmetros são carregados diretamente no estado
inicial, variáveis externas à função são invisíveis dentro dela,
implementando o escopo de função isolado exigido pela regra (T-FuncDecl).

\section{Conclusão}

Neste capítulo, mostramos que o verificador de tipos é uma
transcrição quase literal do sistema de tipos, e que a
correspondência entre regras e código é precisa e sistemática.

Para TExp, a ausência de variáveis permite que o verificador seja uma
função pura sem contexto externo. A mônada
\haskell{ExceptT String Identity} é suficiente.

Para TLine, a introdução de variáveis exige um contexto de tipos
\(\Gamma\) mantido como estado da mônada. O verificador deve
consultar o contexto para variáveis e estendê-lo para declarações, o
que é naturalmente capturado por
\haskell{StateT Ctx (ExceptT String Identity)}.

Para TWhile, a adição de blocos com escopo local exige que o
verificador salve e restaure o contexto ao redor de cada bloco. As
regras (T-If) e (T-While) formalizam essa exigência ao manter o
contexto de saída igual ao de entrada.

Para TImp, funções e registros introduzem dois novos ambientes
globais imutáveis (\(\Delta\) e \(\Phi\)) e um contexto de retorno
por função (\(\rho_{\mathit{ret}}\)), todos naturalmente capturados
por \haskell{ReaderT TcEnv}. O contexto de variáveis \(\Gamma\)
permanece no \haskell{StateT}, crescendo com declarações e sendo
restaurado ao sair de cada bloco ou função.

O padrão que emerge é geral: a estrutura da mônada de verificação
espelha a estrutura dos ambientes semânticos: informação imutável
vai em \haskell{ReaderT}, informação mutável vai em
\haskell{StateT}, e falhas vão em \haskell{ExceptT}. O sistema de
tipos é a especificação; o verificador de tipos é a sua implementação
direta.
